\clearpage
\setcounter{chapter}{8}
\chaptermark{光在晶体中的传播}
\pdfbookmark[0]{第八章 光在晶体中的传播}{chapter8}
\problem{8-1}{
一束线偏振的钠黄光正入射于一方解石晶体，其振动方向与晶体主截面之夹角为 $20^\circ$。试求传播于晶体中的 $\mathrm o,\mathrm e$ 两束光的相对振幅和相对强度。此时取 $n_{\mathrm o}=1.658$，$n_{\mathrm e}=1.486$。
}
\problem{8-2}{
参见题图（a），有两个相同的冰洲石晶体 $C_1$ 和 $C_2$ 前后放置，而两者主截面之夹角为 $\alpha$，图（a）中显示的为 $\alpha=0$ 的样子。一强度为 $I_0$ 的一束自然光正入射于 $C_1$，试分别求出 $\alpha=0^\circ$、$45^\circ$、$90^\circ$ 和 $180^\circ$ 四种情形下，从 $C_2$ 出射之光的强度和数目。这里忽略反射和吸收等耗损。

提示：先在纸面上画好光轴 $\mathrm e_1,\mathrm e_2$ 取向和光振幅矢量的投影如图（b）。
\tikzfig{fig_8_2}{8-2}
}
\problem{8-3}{
一水晶薄片厚 $0.850\,\mathrm{mm}$，其光轴平行于表面。现用一绿光束 $546.1\,\mathrm{nm}$ 正入射于这水晶片；已知水晶对波长为 $546.1\,\mathrm{nm}$ 的绿光的主折射率为 $n_{\mathrm o}=1.5462$，$n_{\mathrm e}=1.5554$。求：
\begin{enumerate}
\item $\mathrm o,\mathrm e$ 两光束在晶片中的光程。
\item 两者经晶片后的相位差。
\end{enumerate}
}
\problem{8-4}{
一束钠黄光以 $50^\circ$ 入射角射向一冰洲石平板，其光轴垂直于入射面，求晶体板中 $\mathrm o$ 光束与 $\mathrm e$ 光束之夹角。已知此时的两个主折射率为 $n_{\mathrm o}\approx1.658$，$n_{\mathrm e}\approx1.486$。
}
\problem{8-5}{
一束钠黄光掠入射于一块冰的平板上，其光轴垂直于入射面。设冰板厚度 $d$ 为 $4.2\,\mathrm{cm}$，求 $\mathrm o$ 光束和 $\mathrm e$ 光束射到平板对面上两点之间隔 $\Delta x$。已知此时的两个主折射率为 $n_{\mathrm o}=1.3090$，$n_{\mathrm e}=1.3104$。
}
\problem{8-6}{
用 ADP 晶体制成 $50^\circ$ 顶角的棱镜，其光轴平行于折射棱即垂直于棱镜的主截面，主折射率为 $n_{\mathrm o}=1.5246$，$n_{\mathrm e}=1.4792$。试求 $\mathrm o$ 光和 $\mathrm e$ 光的最小偏向角及二者之差。
}
\problem{8-7}{
一水晶棱镜的顶角为 $60^\circ$，其光轴垂直于棱镜的主截面，一束钠黄光以近似满足最小偏向角的方式入射于这棱镜。$n_{\mathrm o}=1.54425$，$n_{\mathrm e}=1.55336$。现用焦距为 $1\,\mathrm m$ 的透镜聚焦，试求 $\mathrm o$ 光焦点与 $\mathrm e$ 光焦点之间隔 $\Delta x$。
}
\problem{8-8}{
试针对下列四种情形，定性地画出在晶体内、外 $\mathrm e$ 光的射线方向和偏振方向及其波法线方向，图中虚线表示单轴晶体的光轴方向。
\tikzfig{fig_8_8}{8-8}
}
\problem{8-9}{
方解石晶体对于汞绿光的主折射率为 $n_{\mathrm o}=1.66168$，$n_{\mathrm e}=1.48792$，试问在这晶体内部绿光的波法线与其射线之最大夹角 $\alpha_{\mathrm M}$ 为多少？此时波法线与光轴之夹角 $\theta_0$ 为多少？此时射线与光轴之夹角 $\xi_0$ 为多少？
}
\problem{8-10}{
一厚度 $d$ 为 $20\,\mathrm{mm}$ 的冰洲石晶体，其光轴与表面之夹角为 $40^\circ$。一束钠黄光正入射于该晶体。试求晶体内 $\mathrm e$ 光传播方向，和出射双光束的位移量 $\Delta$。要求画图示意。

提示：参考正文 8.2 节例题 2 和图 8.10。
}
\problem{8-11}{
根据习题 8-10 题所获悉的数据，进一步求出晶体内该 $\mathrm e$ 光的射线折射率 $n_{\mathrm r}$ 和法线折射率 $n_{\mathrm N}$，以及相应的光程 $L_{\mathrm r}$ 和 $L_{\mathrm N}$。
}
\problem{8-12}{
如图所示，光线入射角为 $i$，光轴倾角为 $\beta$，主截面与入射面重合，且
\[
i=\frac{\pi}{2}-|\beta|,
\]
即晶体外光线的入射方向与晶体内的光轴方向一致。
\begin{enumerate}
\item 你认为此种情形下光线在晶体内会发生双折射吗？
\item 如是，试求出 $\mathrm o$ 光、$\mathrm e$ 光的射线方向。设该晶体的两个主折射率为 $n_{\mathrm o}=1.88$，$n_{\mathrm e}=1.44$；$i=30^\circ$。
\end{enumerate}
提示：参考 8.2 节例题 4。
\tikzfig{fig_8_12}{8-12}
}
\problem{8-13}{
一个沃拉斯顿棱镜顶角 $\alpha=15^\circ$，求两条出射光线之夹角。
}
\problem{8-14}{
用方解石或石英制作针对钠黄光的 $\lambda/4$ 波晶片。
\begin{enumerate}
\item 试求它们的最小厚度 $d_1,d_2$ 各为多少？
\item 若用汞绿光入射 $d_1$ 片或 $d_2$ 片，其相位延迟 $\delta_1,\delta_2$ 各为多少？
\end{enumerate}
}
\problem{8-15}{
一方解石棱镜其顶角为 $30^\circ$，光轴垂直棱镜主截面，一束钠黄光从左侧正入射。试从这棱镜右侧出射两束光的射线方向和偏振方向。要求作出示意图。

注：棱镜的主截面定义为与棱镜的折射棱边正交的那个截面。
}
\problem{8-16}{
一单色线偏振的窄光束正入射于一方解石晶体而出现双折射，其偏振方向与晶体主截面之夹角为 $30^\circ$；一尼科耳棱镜或一偏振片置于晶体后面，其主截面与原入射线偏振成 $50^\circ$ 角。求最终出射的两束光之光强比值。
}
\problem{8-17}{
一块水晶片厚度为 $0.850\,\mathrm{mm}$，光轴平行于表面，一波长为 $5461\,\text{\AA}$ 的绿光正入射这水晶片。求 $\mathrm o,\mathrm e$ 两光束通过这晶片的光程差和相位差。
}
\problem{8-18}{
现有三块自然云母片，经螺旋测厚仪确定其厚度分别为 $0.710\,\mathrm{mm}$、$0.824\,\mathrm{mm}$ 和 $0.938\,\mathrm{mm}$。针对钠黄光，试问：
\begin{enumerate}
\item 选用其中哪一片作为 $\lambda/2$ 波晶片最为合适？
\item 其中是否有一片较合适于作为 $\lambda/4$ 波晶片？
\item 要求对上述被选中的波晶片的偏差作出估量，并以相位差的度数明示。
\end{enumerate}
提示：参考 8.3 节例题 2——剥离云母片的合适厚度。
}
\problem{8-19}{
让一束椭圆偏振光，先后通过一 $\lambda/4$ 片和一张偏振片 $P$，旨在监测椭圆光的特征。在转动 $P$ 过程中出现了消光，此时 $\lambda/4$ 片的光轴与 $P$ 片的透振方向之夹角为 $22^\circ$。
\begin{enumerate}
\item 求入射的椭圆光之长短轴之比值。
\item 是否可以凭借这个实验对入射椭圆光的左右旋性作出判断。设 $\lambda/4$ 片提供的附加相位差 $\delta'_{\mathrm{oe}}=+\pi/2$。
\end{enumerate}
}
\problem{8-20}{
一强度为 $I_0$ 的右旋圆偏振光正入射于一 $\lambda/4$ 片，然后再通过一偏振片 $P$，其透振方向相对 $\lambda/4$ 片光轴方向顺时针旋 $15^\circ$。
\begin{enumerate}
\item 求最后出射光强。设 $\lambda/4$ 片提供的附加相位差 $\delta'_{\mathrm{oe}}=-\pi/2$。
\item 若 $P$ 之透振方向逆时针旋 $15^\circ$，出射光强为多少？
\end{enumerate}
}
\problem{8-21}{
凭借一波晶片是否可能将一椭圆偏振光改变为一圆偏振光？下面的习题有助于解答这个问题。设一右旋椭圆偏振光其长短轴之比值为 $A_y/A_x=3$，正入射于一波晶片，其光轴方向相对于 $y$ 轴或 $x$ 轴为 $45^\circ$，于是获得 $\mathrm o,\mathrm e$ 两个正交振动，且两者等振幅，$A_{\mathrm o}=A_{\mathrm e}$。进一步讨论相位差问题：
\begin{enumerate}
\item 对波晶片的入射点而言，$\mathrm o,\mathrm e$ 两个振动的相位差 $\delta_{\text{入}}$ 为多少？
\item 波晶片提供的附加相位差 $\delta_{\mathrm{oe}}$ 应满足何种条件可使出射光为圆偏振光？
\end{enumerate}
}
\problem{8-22}{
一对偏振片 $P_1$ 和 $P_2$ 其透振方向之夹角为 $30^\circ$，其间插入一波晶片，其光轴方向恰在那 $30^\circ$ 角的平分线上。设入射光为一单色自然光，振幅为 $A$，光强为 $I_0$，忽略反射、吸收等损耗。（要求作出示意图）
\begin{enumerate}
\item 求从波晶片出射的 $\mathrm o$ 光、$\mathrm e$ 光的振幅和光强。
\item 求投影于 $P_2$ 透振方向的那两个振动成分的振幅和光强。
\item 求最终通过 $P_2$ 的输出光强 $I_2$。设波晶片为 $\lambda/4$ 片。
\end{enumerate}
}
\problem{8-23}{
在一对正交的偏振片 $P_1$ 和 $P_2$ 之间插入一块 $\lambda/4$ 片，其光轴与 $P_1$ 透振方向成 $60^\circ$ 角。一强度为 $I_0$ 的单色自然光正入射于该系统，求出射光的强度。忽略反射、吸收损耗，要求作图示意。
}
\problem{8-24}{
一块波晶片由方解石切割而成，其厚度 $d=25\,\mu\mathrm m$，现将其置于一对正交偏振片之间，波晶片的光轴方向恰巧与透振方向成 $45^\circ$ 角。设在可见光波段，方解石的两个主折射率的平均值为 $\bar n_{\mathrm o}\approx1.6678$，$\bar n_{\mathrm e}\approx1.4904$。要求作图示意。
\begin{enumerate}
\item 在可见光谱区哪些波长的光几乎不能通过该系统？
\item 若将第二个偏振片转过 $90^\circ$，又是哪些波长的光几乎不能通过这系统？
\end{enumerate}
}
\problem{8-25}{
一右旋椭圆偏振光相继通过一块波晶片和一块偏振片。波晶片由负晶体制成，它对入射光的有效光程差为 $(n_{\mathrm o}-n_{\mathrm e})d=\lambda/6$，且其光轴方向已对准椭圆光的短轴方向；偏振片透振方向沿光轴左旋 $30^\circ$ 角。入射光总光强为 $I_0$，其极大、极小光强比为 $4$。
\begin{enumerate}
\item 求出射光强 $I$。
\item 若入射光改变为左旋，而其他条件不变，出射光强为多少？
\item 在偏振片转动过程中，出射光强的极大值和极小值分别为多少？
\end{enumerate}
}
\problem{8-26}{
一左旋圆偏振光相继通过一波晶片和一偏振片。波晶片由正晶体制成，其对入射光产生的有效光程差为 $(n_{\mathrm e}-n_{\mathrm o})d=\lambda/3$；偏振片透振方向沿光轴右旋 $30^\circ$ 角。设入射光强度为 $I_0$。
\begin{enumerate}
\item 参与偏振片透振方向相干叠加的两个振动的相位差 $\delta_{\mathrm{oe}}$ 为多少（$\mathrm{rad}$）？
\item 求出射光强。要求作出示意图。
\item 若入射光改为右旋，而其他条件不变，求出射光强。
\end{enumerate}
}
\problem{8-27}{
汞灯的 $4047\,\text{\AA}$ 紫色平行光束正入射于一偏振光干涉系统，该系统即是在两个正交偏振片 $P_1$ 和 $P_2$ 之间放置一楔形水晶薄棱镜，棱镜顶角 $\alpha=0.5^\circ$，光轴平行棱边且与偏振片透振方向成 $45^\circ$ 角。
\begin{enumerate}
\item 通过 $P_2$ 看到的干涉图样是怎样的？（要求作图示意。）
\item 相邻暗纹的间隔 $d$ 为多少？已知 $n_{\mathrm o}=1.5572$，$n_{\mathrm e}=1.5667$。
\item 若将 $P_2$ 转过 $90^\circ$，干涉图样有何变化？
\item 维持 $P_1,P_2$ 正交，但将水晶棱镜的光轴方向转过 $45^\circ$，使之与 $P_1$ 透振方向平行，干涉图样有何变化？
\end{enumerate}
}
\problem{8-28}{
一束钠黄光正入射于一巴比涅补偿器（参见正文图 8.22（b））。该补偿器由水晶制成，其楔角 $\alpha=2.75^\circ$，并置于两个正交偏振片之间。已知对钠黄光水晶的两个主折射率为 $n_{\mathrm o}=1.54425$，$n_{\mathrm e}=1.55336$。
\begin{enumerate}
\item 从补偿器出射的两束平行光之夹角 $\Delta\theta$ 为多少？
\item 通过 $P_2$ 而出现的干涉条纹的间隔 $\Delta x$ 为多少？
\end{enumerate}
}
\problem{8-29}{
一巴比涅补偿器置于两个正交偏振片 $P_1$ 和 $P_2$ 之间，且其两个透振方向与补偿器的两个光轴方向互成 $45^\circ$。通过 $P_2$ 观察到中央有一条暗线且暗纹间隔为 $a$。今以同波长的椭圆偏振光直接照射这块补偿器以替代 $P_1$，发现这条暗纹移动 $b$ 距离。
\begin{enumerate}
\item 求入射椭圆光相对补偿器两个正交光轴所构成的坐标架，其两个振动之相位差 $\delta_0$ 为多少？
\item 若椭圆光的长短轴恰巧对准了补偿器的两个正交的光轴，此时 $a,b$ 有何关系？
\item 若此时 $P_2$ 透振方向与补偿器中一光轴之夹角为 $\theta$，试给出 $\theta$ 角与入射椭圆光的长短轴比值之关系。
\end{enumerate}
}
\problem{8-30}{
已知石英对钠黄光的旋光率 $\alpha=21.72\,({}^\circ)/\mathrm{mm}$，试求左、右旋圆偏振光传播于石英晶体沿光轴方向的折射率之差 $\Delta n$。
}
\problem{8-31}{
在两块正交偏振片之间插入一石英旋光晶片，以消除对人眼最敏感的黄绿光 $\lambda=550\,\mathrm{nm}$，设石英对这一波长的旋光率 $\alpha=24\,({}^\circ)/\mathrm{mm}$。
\begin{enumerate}
\item 求该晶片的最小厚度 $d_{\mathrm m}$。
\item 若转动其中一偏振片使两个透振方向平行，此时晶片最小厚度 $d'_{\mathrm m}$ 应当为多少？
\end{enumerate}
}
\problem{8-32}{
一石英棒之长度为 $5.639\,\mathrm{cm}$，其端面垂直于光轴，被置于两个正交偏振片之间。现有一束白光正入射于此棒的端面，并用光谱仪观察该系统的透射光。
\begin{enumerate}
\item 试在一张坐标纸上绘制偏振面转角随波长而变化的曲线 $\varphi(\lambda)$，其波长范围限于可见光谱区（$400$--$760\,\mathrm{nm}$），所需石英旋光率数据 $\alpha(\lambda)$ 可查正文表 8.5。
\item 从 $\varphi(\lambda)$ 曲线中可以找出一系列离散的波长，它们将在光谱仪中消失。试分别给出紫端和红端前三个这种波长。
\end{enumerate}
}
\problem{8-33}{
一块石英片其表面垂直于光轴，恰好抵消了 $10\,\mathrm{cm}$ 长、浓度为 $0.20\,\mathrm{g/cm^3}$ 的麦芽糖溶液对钠黄光所产生的旋光效应。试求这石英片的厚度。已知石英对此波长的旋光率 $\alpha=21.75\,({}^\circ)/\mathrm{mm}$，麦芽糖的旋光率 $[\alpha]=144\,({}^\circ)/(\mathrm{dm}\cdot\mathrm g\cdot\mathrm{cm}^{-3})$。
}
\problem{8-34}{
一长度为 $15\,\mathrm{cm}$ 的左旋葡萄糖溶液，使钠黄光的偏振面旋转了 $25.6^\circ$。已知葡萄糖溶液的旋光率 $[\alpha]=-51.4\,({}^\circ)/(\mathrm{dm}\cdot\mathrm g\cdot\mathrm{cm}^{-3})$，求该溶液的葡萄糖浓度。
}
\problem{8-35}{
将 $14.5\,\mathrm g$ 的蔗糖块溶于水而获得一体积为 $50\,\mathrm{cm^3}$ 的溶液，并灌入长度为 $15\,\mathrm{cm}$ 的量糖计中，测出钠光束偏振面向右旋转了 $16.8^\circ$。求这蔗糖样品中所含非旋光性杂质的百分比。已知纯糖溶液的旋光率 $[\alpha]=66.5\,({}^\circ)/(\mathrm{dm}\cdot\mathrm g\cdot\mathrm{cm}^{-3})$。
}
\problem{8-36}{
已知右旋石英对 $\lambda=762\,\mathrm{nm}$ 的左、右旋圆偏振光的折射率分别为 $n_{\mathrm L}=1.53920$，$n_{\mathrm R}=1.53914$，试求石英对该波长的旋光率。
}
\problem{8-37}{
（接上题）用该石英制成一晶体棒，其端面垂直于光轴，一束波长为 $762\,\mathrm{nm}$ 的线偏振光正入射于这晶体棒。求沿光轴方向相距 $2.000\,\mathrm{mm}$ 的两点，其偏振面之夹角 $\Delta\psi$ 为多少？其振动之有效相位差 $\Delta\varphi$ 为多少？
}
\problem{8-38}{
用磁感应强度 $B=0.6$ 特斯拉的磁场加于一块长 $10\,\mathrm{cm}$ 的轻火石玻璃上，求偏振面旋转角（度数）。已知该材料的韦尔代常数 $V=9.22\,\mathrm{rad/(T\cdot m)}$。
}
\problem{8-39}{
用 $20\,\mathrm{cm}$ 长的某种液体观察法拉第效应，若加上 $0.8\,\mathrm T$ 的磁场，观测到线偏振光的振动面旋转了 $65^\circ$，求该液体的韦尔代常数 $V$。
}
\problem{8-40}{
用长度 $5\,\mathrm{cm}$ 的重火石玻璃棒作为光学隔离器的元件，求所施加的磁场应当为多大。已知其韦尔代常数为 $30\,\mathrm{rad/(T\cdot m)}$。
}
\problem{8-41}{
用克尔常数 $K=2.44\times10^{-12}\,\mathrm{m/V^2}$ 的硝基苯液体制成克尔盒，其极板长 $2.8\,\mathrm{cm}$，两板间距 $0.6\,\mathrm{cm}$。
\begin{enumerate}
\item 若要求克尔盒装置有最大光强输出，则应当施加多大电压？
\item 这一最大输出光强是输入自然光强度的百分之几？
\end{enumerate}
}
\problem{8-42}{
将硝基苯注入克尔盒，其板长 $3.0\,\mathrm{cm}$，两极板间距 $0.75\,\mathrm{cm}$，所加电压为 $22\,\mathrm{kV}$。
\begin{enumerate}
\item 求从克尔盒出射的两个正交振动之间的相位差 $\delta'$，要求作图示意。
\item 若光强为 $I_0$ 的自然光束入射于这一系统，求最终输出光强为多少？
\end{enumerate}
}
